% paper/sections/11b6_dccd_dissipation.tex
%
% §11.3.6  DCCD フィルタ：分散・散逸特性と選択的高周波減衰の評価
%          実験スクリプト: experiments/dccd_comparison.py
%          結果データ:    results/dccd_comparison/

\paragraph{主張．}

本実験が示す主命題は以下の二点である：

\begin{center}
\textbf{%
  DCCD の 10 次選択フィルタ（§~\ref{sec:advection_dccd_design}）は，%
  不連続プロファイルに対して CCD の過大な変動（全変動：TV）を大幅に抑制しつつ，%
  滑らかなプロファイルでは CCD と同等の高精度を維持する（選択的高周波散逸）．%
}
\end{center}

CLS 法では界面 $\psi$ が $0$–$1$ 遷移層を持つ準不連続プロファイルである．
CCD（フィルタなし）は高波数成分に対して分散誤差が蓄積し，
界面近傍に非物理的な振動（Gibbs 型）を生じやすい．
DCCD フィルタはスペクトル領域で高波数成分のみを選択的に減衰させるため，
滑らかな領域の精度を犠牲にせずに振動を抑制する（§~\ref{sec:advection_dccd_design}）．

\paragraph{実験設定．}

1次元線形移流方程式
\begin{equation}
  \frac{\partial u}{\partial t} + c\,\frac{\partial u}{\partial x} = 0,
  \qquad x \in [0,1),\quad c = 1,\quad T = 1
  \label{eq:linear_advection_dccd}
\end{equation}
を $N=256$ の周期格子，CFL $=0.4$，RK4 で 1 周期（$T=1$）移流する．
初期条件として 3 種類のプロファイルを用いる：
\begin{itemize}[noitemsep]
  \item \textbf{Square}：不連続矩形パルス（$x \in [0.3, 0.5)$ で $u=1$，他は $0$）
  \item \textbf{Triangle}：$C^0$ 三角波（頂点 $x=0.4$ で $u=1$，$[0.3,0.5]$ 台）
  \item \textbf{Smooth tanh}：双曲線正接型平滑化矩形（準不連続・滑らか）
\end{itemize}
比較スキーム：O2（2次中心差分），O4（4次中心差分），CCD（6次），
DCCD（CCD $+$ 10次選択フィルタ，$\alpha_f=0.4$），WENO5（5次重み付き ENO）．

評価指標：
$L_2 = \sqrt{\text{mean}(u-u_\text{exact})^2}$，
$\text{TV} = \sum_i |u_{i+1}-u_i|$（周期ラップ込み）．

\paragraph{結果．}

\begin{table}[htbp]
\centering
\caption{1周期後の $L_2$ 誤差（$N=256$，CFL$=0.4$，太字：各初期条件での最良値）．
  スクリプト: \texttt{experiments/dccd\_comparison.py}．}
\label{tab:dccd_l2_bench}
\begin{tabular}{lrrrrr}
\toprule
初期条件 & O2 & O4 & CCD & DCCD & WENO5 \\
\midrule
Square（不連続） & $1.41\times10^{-1}$ & $9.08\times10^{-2}$ & $4.80\times10^{-2}$ & $\mathbf{4.23\times10^{-2}}$ & $6.35\times10^{-2}$ \\
Triangle（$C^0$） & $2.41\times10^{-2}$ & $5.84\times10^{-3}$ & $\mathbf{1.37\times10^{-3}}$ & $1.41\times10^{-3}$ & $3.90\times10^{-3}$ \\
Smooth tanh & $3.75\times10^{-2}$ & $2.46\times10^{-3}$ & $\mathbf{2.57\times10^{-5}}$ & $3.75\times10^{-5}$ & $9.66\times10^{-4}$ \\
\bottomrule
\end{tabular}
\end{table}

\begin{table}[htbp]
\centering
\caption{1周期後の全変動（TV）．
  Exact 値は初期条件の TV（Square: 2.000, Triangle: 1.969, tanh: 2.000）．
  太字：各初期条件での最良値（TV $\to$ exact に最近）．}
\label{tab:dccd_tv_bench}
\begin{tabular}{lrrrrr}
\toprule
初期条件 & O2 & O4 & CCD & DCCD & WENO5 \\
\midrule
Square（不連続） & $18.40$ & $17.17$ & $10.83$ & $3.15$ & $\mathbf{2.00}$ \\
Triangle（$C^0$） & $2.60$ & $2.37$ & $2.09$ & $1.97$ & $\mathbf{1.91}$ \\
Smooth tanh & $2.82$ & $2.02$ & $2.00$ & $2.00$ & $\mathbf{2.00}$ \\
\bottomrule
\end{tabular}
\end{table}

\begin{figure}[htbp]
  \centering
  \includegraphics[width=0.95\linewidth]{%
    ../results/dccd_comparison/figures/04_waveforms_panel.png}
  \caption{%
    3 種の初期条件で $T=1$ 後の波形比較（O2・O4・CCD・DCCD・WENO5）．
    左から Square, Triangle, tanh．
    Square では CCD に顕著な Gibbs 型振動が残るが，DCCD フィルタにより大幅に抑制される．
    DCCD の TV は $3.15$（CCD の $10.83$ に対して約 1/3）であり，WENO5 の $2.00$ に近い．
    Smooth tanh では CCD と DCCD の波形はほぼ重なり，フィルタが滑らかな成分に影響しないことを示す
    （$L_2$: CCD $2.57\times10^{-5}$ vs DCCD $3.75\times10^{-5}$）．}
  \label{fig:dccd_waveforms_bench}
\end{figure}

\begin{tcolorbox}[resultbox, title={DCCD フィルタの選択的散逸効果}]
\begin{itemize}[noitemsep]
  \item \textbf{不連続（Square）：} DCCD は CCD に比べ TV を $10.83 \to 3.15$ に低減（約 $1/3$）し，
        Gibbs 型振動を大幅に抑制する．$L_2$ 誤差も最良（$4.23\times10^{-2}$）．
        WENO5 が TVD 条件を満たすため最良 TV（$2.00$）だが，$L_2$ は DCCD より悪い（$6.35\times10^{-2}$）．
  \item \textbf{平滑（Smooth tanh）：} DCCD と CCD の $L_2$ 誤差は同オーダー
        （$3.75\times10^{-5}$ vs $2.57\times10^{-5}$，1.5 倍以内），TV も同値（$2.000$）．
        フィルタが高波数成分のみを対象とし，滑らかな成分を保護することを示す．
  \item \textbf{CLS 法への含意：} CLS の界面 $\psi$ は Smooth tanh に相当する平滑化プロファイルであるため，
        DCCD は CCD と同等の精度で界面を移流しつつ，界面近傍のわずかな非物理的振動を抑制する．
        これは §~\ref{sec:interface_advection_bench}（Zalesak・単一渦）の低体積誤差の理論的根拠となる．
\end{itemize}
\end{tcolorbox}
